% Figure: the denoiser drawn as a stack of layers, and the LoRA adapter on one linear layer.
% Facts come from the published model repository and model card.
% Requires in the preamble: \usepackage{tikz} + \usetikzlibrary{positioning,arrows.meta},
%                            and amsmath for \frac in the adapter formula.
\begin{figure}[t]
  \centering
  \begin{tikzpicture}[
    x=1cm, y=1cm,
    font=\small,
    block/.style={rectangle, rounded corners=2pt, draw=black!70, line width=0.5pt,
      fill=blue!10, align=center, inner sep=4pt, minimum height=9mm},
    inputblock/.style={block, fill=green!12},
    stageblock/.style={block, fill=orange!15},
    noteblock/.style={rectangle, rounded corners=2pt, draw=black!45, dashed,
      fill=red!8, align=left, inner sep=3pt, font=\scriptsize},
    flow/.style={->, >=Latex, line width=0.7pt, draw=black!75},
    feed/.style={->, >=Latex, line width=0.7pt, draw=black!75, dashed},
    arrowlab/.style={font=\scriptsize, fill=white, inner sep=1pt},
  ]

  % ================= left part: the denoiser as a stack of layers =================
  \node[inputblock, anchor=north, text width=6.40cm, minimum height=0.90cm] (lat) at (6.54,10.55)
    {\textbf{latent token}, 128 channels};
  \node[block, text width=6.40cm, minimum height=0.90cm, below=0.32cm of lat] (proj)
    {\textbf{input projection} to width 3072};
  \node[block, text width=6.40cm, minimum height=1.15cm, below=0.32cm of proj] (tstep)
    {\textbf{timestep embedding}\\[1pt]\scriptsize sets scale and shift for every block};
  \node[block, text width=6.40cm, minimum height=1.15cm, below=0.32cm of tstep] (ds)
    {\textbf{5 double-stream blocks}\\[1pt]\scriptsize text and image tokens stay apart,
     joint attention};
  \node[block, text width=6.40cm, minimum height=0.90cm, below=0.32cm of ds] (cat)
    {\textbf{concatenate the two streams}};
  \node[block, text width=6.40cm, minimum height=1.15cm, below=0.32cm of cat] (ss)
    {\textbf{20 single-stream blocks}\\[1pt]\scriptsize one stream, joint attention};
  \node[block, text width=6.40cm, minimum height=1.15cm, below=0.32cm of ss] (fin)
    {\textbf{final layer}\\[1pt]\scriptsize adaptive layer norm, project to 128 channels};
  \node[inputblock, text width=6.40cm, minimum height=0.90cm, below=0.32cm of fin] (latout)
    {\textbf{latent}, to the decoder};
  \foreach \a/\b in {lat/proj, proj/tstep, tstep/ds, ds/cat, cat/ss, ss/fin, fin/latout}
    {\draw[flow] (\a.south) -- (\b.north);}

  % the text encoder enters through joint attention (dashed = conditioning path)
  \node[block, text width=2.10cm, left=0.70cm of cat] (enc)
    {\textbf{Qwen3-4B}\\[1pt]\scriptsize text encoder};
  \coordinate (txtcorner) at (enc.north |- ds.center);
  \draw[feed] (enc.north) -- (txtcorner) -- (ds.west);
  \path (txtcorner) -- (ds.west) coordinate[midway] (txtmid);
  \node[arrowlab, align=center, text width=1.70cm, anchor=south] at ([yshift=0.14cm]txtmid)
    {text tokens,\\width 7680};

  % the timestep embedding modulates every block (dashed = conditioning path)
  \draw[feed] (tstep.east) -- ++(0.35,0) coordinate (tbus);
  \draw[draw=black!75, line width=0.7pt, dashed] (tbus) -- (tbus |- fin.center);
  \foreach \s in {ds,ss,fin}{\draw[feed] (tbus |- \s.center) -- (\s.east);}

  % ================= right part: how the adapter attaches =================
  \node[inputblock, text width=2.40cm, minimum height=0.75cm] (layin) at (12.45,10.10)
    {layer input $x$};
  \node[block, text width=1.70cm, minimum height=2.35cm] (wgt) at (11.55,7.575)
    {\textbf{frozen weight $W$}};
  \node[block, text width=0.50cm, minimum height=0.90cm] (bA) at (13.35,8.30) {\textbf{A}};
  \node[block, text width=0.50cm, minimum height=0.90cm] (bB) at (13.35,6.85) {\textbf{B}};
  \draw[flow] (bA.south) -- (bB.north);

  \coordinate (br) at (12.45,9.30);
  \draw[draw=black!75, line width=0.7pt] (layin.south) -- (br);
  \draw[draw=black!75, line width=0.7pt] (br -| wgt.north) -- (br -| bA.north);
  \draw[flow] (br -| wgt.north) -- (wgt.north);
  \draw[flow] (br -| bA.north) -- (bA.north);

  \coordinate (mg) at (12.45,6.05);
  \draw[flow] (wgt.south) -- (mg -| wgt.south);
  \draw[flow] (bB.south) -- (mg -| bB.south);
  \draw[draw=black!75, line width=0.7pt] (mg -| wgt.south) -- (mg -| bB.south);
  \node[fill=white, inner sep=0.5pt] at (mg) {$+$};
  \node[inputblock, text width=1.85cm, minimum height=0.75cm] (layout) at (12.45,5.30)
    {output};
  \draw[flow] (mg) -- (layout.north);
  \node[font=\small] (formula) at (12.45,4.45) {$Wx + \frac{\alpha}{r}BAx$};

  \node[font=\scriptsize, align=left, text width=1.05cm, anchor=west] (abl) at (13.90,7.575)
    {rank $r$\\trained};

  % ================= notes =================
  \node[noteblock, text width=2.10cm, anchor=north west] at (0.12,3.40)
    {This layer stack is the published layout. We did not change it.};
  \node[noteblock, text width=4.40cm, anchor=north west] at (10.60,3.95)
    {Only A and B are trained. $W$ never changes, so the adapter is small and it can be
     removed again.};
  \node[noteblock, text width=4.40cm, anchor=north west] at (10.60,2.10)
    {The adapter has 100 modules. It teaches the base model the head swap task, so it changes
     what the model does, not how it looks.};

  \end{tikzpicture}
  \caption{The denoiser as a stack of layers. The text encoder and the timestep embedding enter
    as conditioning, and the adapter adds a small trained branch beside a frozen weight.}
  \label{fig:architecture_lora}
\end{figure}
